\documentclass[11pt,a4paper]{article}
\usepackage[margin=25mm]{geometry}
\usepackage{amsmath,amssymb,amsthm,hyperref}
\setlength{\emergencystretch}{2em}
\newtheorem{lemma}{Lemma}
\newtheorem{theorem}[lemma]{Theorem}
\newtheorem{corollary}[lemma]{Corollary}
\newcommand{\gen}[1]{\langle#1\rangle}
\title{Deleting an anchor from a paired string action}
\author{Human Compression project}\date{}
\begin{document}\maketitle
Let $G=\gen{h_0,\ldots,h_{r-1}}$ be transitive on $2r+1$ points.
The marks are nonidentity involutions satisfying the string relations.
An edge of colour $i$ is eligible if its endpoints belong to different
orbits of $G_i=\gen{h_j:j\ne i}$. Assume that each colour has at least
two eligible edges. A first anchor is moved by $h_0$ and fixed by all
later marks; a last anchor is defined in the reversed order.

We acquire the complete generic paired-cut, equality-connectedness,
square-neighbour and ladder proofs from the L044 recovery. In particular,
in degree $2r>8$ their actual equality ladder has no first or last
anchor: every point moved by an end mark is moved by another mark.
We also acquire the ambient interval deficit. If an intransitive
contiguous subfamily has $s$ marks, $f$ fixed points and $c$ nontrivial
orbits, then
\[
                         2s\le 2r+1-f-c.                 \tag{1}
\]
These statements use neither parity nor the intersection property.

\begin{lemma}\label{one}
There is at most one first anchor. If it exists and is $b$, then
$G_0$ has exactly three orbits: $\{b\}$ and two nontrivial orbits.
\end{lemma}
\begin{proof}
Suppose two first anchors $b,w$ exist. For $H=G_0$, (1) gives
$f+c\le3$, while $f\ge2$ and $c\ge1$. Thus $f=2,c=1$.
Let $x=b^{h_0}$ and $z=w^{h_0}$. The involution cannot pair $b,w$,
since that pair would be invariant under every generator. Consequently
$x,z$ are distinct points of the unique nontrivial $H$-orbit $V$,
whose size is $2r-1$. Every $h_i$, $i\ge2$, fixes $x,z$ by
commutation with $h_0$. Transitivity on $V$ makes $h_1$ move both.

The action of $H$ on $V$ is faithful, since $H$ fixes $b,w$.
Every mark remains nonidentity. Each original eligible edge of a
colour $i>0$ lies in $V$ and remains eligible for the smaller
omission subgroup of $H$. We have therefore obtained a paired
string action of degree $2(r-1)+1$, rank $r-1$, and two first
anchors. Induction is impossible at rank one, where an involution
cannot be transitive on three points. This proves the first assertion.

Now let $b$ be a first anchor. The same interval bound gives
$f+c\le3$ for $G_0$, with $f,c\ge1$. If $f=c=1$, the only
eligible $h_0$-edge is its edge at $b$, contrary to the paired
hypothesis. The possibility $f=2,c=1$ would give a second first
anchor: any further $G_0$ fixed point must be moved by $h_0$,
by full transitivity. It has just been excluded. Hence $f=1,c=2$.
\end{proof}

\begin{lemma}\label{delete}
Suppose $r\ge5$ and there are anchors at both ends. Each end
colour has exactly two eligible edges.
\end{lemma}
\begin{proof}
Write $h_0=(b\ x)\eta$, where $b$ is the first anchor. Delete
$b$ and replace $h_0$ by $\eta$. The remaining permutation graph
is connected, since $b$ was a pendant vertex. The string relations
persist: for $i\ge2$, commutation with $h_0$ and fixation of $b$
give fixation of $x$ and commutation with $\eta$. The latter mark
is nonidentity because a second eligible $h_0$-edge existed.

Every omission orbit on the remaining points is unchanged. For
an omission containing $h_0$, deleting its pendant vertex $b$
does not change paths between other vertices. In the omission
of $h_0$, that vertex was already isolated. Thus every later
colour retains its eligible edges, and colour zero loses exactly
the edge $bx$. The last anchor survives: it is fixed by $h_0$
and is therefore neither $b$ nor $x$.

If colour zero retained at least two eligible edges, the modified
degree-$2r$ action would satisfy all hypotheses of the acquired
paired equality ladder. That ladder has no last anchor, a
contradiction. Hence exactly one eligible edge remains. Reversing
the labels proves the other endpoint assertion. No intersection
property is assumed for the modified action.
\end{proof}

\begin{theorem}[The four-point head]\label{head}
Suppose $r\ge5$ and there are distinct first and last anchors $b,c$.
There are distinct points $b,x,u,v,z$ such that, with
$V=\Omega\setminus\{b,x,u,v\}$,
\[
 h_0=(b\ x)(u\ v),\qquad h_1=(x\ u)(v\ z)\beta,       \tag{2}
\]
where $\beta$ is supported in $V\setminus\{z\}$.
The group $K=\gen{h_2,\ldots,h_{r-1}}$ fixes the four omitted
points and is transitive on $V$. On $V$ it is a paired string
action of rank $r-2$ and degree $2(r-2)+1$, with first anchor
$z$ and last anchor $c$. If the original marking has the
intersection property, this actual restricted marking does too.
\end{theorem}
\begin{proof}
By Lemma~\ref{one}, $G_0$ has orbits $\{b\},A,B$; label them so
$x=b^{h_0}\in A$. By Lemma~\ref{delete}, there is exactly one
other eligible $h_0$-edge. Full transitivity forces it to join
$u\in A$ to $v\in B$. The four points are distinct. Every
$h_i$, $i\ge2$, fixes $b,x$. It transports the unique eligible
edge between $A,B$ to itself and preserves each of those orbits,
so it also fixes $u,v$.

Apply (1) to $K$. It gives $f+c\le5$. Since $K$ is nontrivial
and fixes the four displayed points, those are exactly its fixed
points and the remaining $V$ is one nontrivial orbit. If $V$
lay in $A$, then $B=\{v\}$, contradicting Lemma~\ref{one}.
Consequently
\[
                        A=\{x,u\},\qquad B=\{v\}\cup V.
\]
The mark $h_1$ swaps $x,u$ and pairs $v$ with some $z\in V$.
Thus it has the form in (2); initially $h_0=(b\ x)(u\ v)\eta$
with $\eta$ supported in $V$. This $\eta$ centralizes the
transitive $K$-action. The last anchor $c$ is fixed by $h_0,h_1$,
so it belongs to $V$ and is fixed by $\eta$. A centralizer of a
transitive action that fixes a point is the identity. Thus $\eta=1$.

For $i\ge3$, fixation of $v$ and commutation with $h_1$ imply
fixation of $z$. Transitivity of $K$ makes $h_2$ move $z$.
The last anchor $c$ remains a last anchor and differs from $z$,
since $h_1$ moves $z$ and fixes $c$. The tail marks fix the
complement of $V$, and their original eligible edges remain
eligible after passing to smaller omission groups. This proves
the claimed paired action. Its faithful identification with the
actual subgroup $K$ also preserves every subgroup intersection.
\end{proof}

\begin{lemma}\label{alternating}
Under Theorem~\ref{head}, $G$ contains $\operatorname{Alt}(\Omega)$.
\end{lemma}
\begin{proof}
In a nontrivial block system, every block has an odd size at least
three. The mark $h_0$ moves only four points and therefore cannot
exchange two blocks. All other marks fix $b$ and preserve its
block. This block is then $G$-invariant, contrary to transitivity.
Thus $G$ is primitive.

Choose $\sigma\in K$ moving $z$ and put $g=h_1\sigma h_1$.
It fixes $b,x,u$ and sends $v$ to some $w\in B\setminus\{v\}$.
The product of $h_0$ and its conjugate by $g$ is a three-cycle
on $u,v,w$ (replacing $g$ by its inverse if required by the
conjugation convention).

For completeness, a primitive group containing a three-cycle
contains the alternating group. Take the supports of all conjugates
of that cycle. Their incidence components are invariant blocks,
so primitivity makes their union connected and equal to $\Omega$.
Three-cycles on intersecting triples generate the alternating group
on their union: intersection in two points gives the usual four-point
alternating group; intersection in one point reduces to that case
by conjugation. Explicitly, for triples $\{1,2,3\}$ and $\{3,4,5\}$,
conjugating the second cycle by the first supplies a triple
$\{1,4,5\}$; these last two triples overlap in two points, and
their four-point alternating group joins the first triple in two
points. Repeatedly adjoining a triple to a connected union now gives
the alternating group on the full component. This proves the assertion.
\end{proof}

\begin{corollary}[The internal factor at the head]\label{beta}
Suppose $r\ge7$ and the original marking has the intersection
property. Then $\beta=1$ in (2). In particular $G$ and its
actual tail $K$ have the same parity: $G$ consists entirely of
even permutations if and only if $K$ does.
\end{corollary}
\begin{proof}
Apply Theorem~\ref{head} to the tail, whose rank $r-2$ is at
least five. In its notation,
\[
 h_2=(z\ p)(q\ w),
\]
and $\gen{h_3,\ldots,h_{r-1}}$ has orbits $\{z\}$, $\{p,q\}$,
and one further orbit containing $c$. The factor $\beta$ fixes
$z,c$ and centralizes this omission group. Its commutation follows
from that of $h_1$ with each $h_i$, $i\ge3$: those marks fix
$x,u,v,z$, so they also commute with the other factor of $h_1$.
The three omission orbits have distinct sizes, and fixation of
$c$ makes $\beta$ the identity on the large orbit. Therefore
$\beta$ is either $1$ or $(p\ q)$.

Suppose $\beta=(p\ q)$. On the five points $b,x,u,v,z$ the
product $h_0h_1$ is a five-cycle; its remaining factor is $\beta$.
Thus $(h_0h_1)^5=\beta$. Conjugating $\beta$ by $h_2$ gives
$(z\ w)$, so
\[
                \delta=(p\ q)(z\ w)\in\gen{h_0,h_1,h_2}.
\]
Lemma~\ref{alternating} applied to the tail puts the same even
permutation $\delta$ in $K$. The original intersection property
then gives $\delta\in\gen{h_2}$, which is false: it is neither
the identity nor $(z\ p)(q\ w)$. Hence $\beta=1$. Both marks
in (2) are now even, proving the parity statement.
\end{proof}

\paragraph{Attribution and remaining boundary.}
The paired selection and ambient orbit-deficit framework comes from
the cited Cameron--Fernandes--Leemans--Mixer sources, with the complete
recovered proofs and disclosed contributions acquired here. Anchor
deletion, the descent excluding two first anchors, and their combination
into (2) are the present deduction. The elementary three-cycle argument
is standard and has been included to make its use explicit.
This component does not acquire the source's near-half classification.
The ranks below five, and the internal factor at rank six, still need
their exact boundary arguments before excluding a symmetric middle
uniformly. No whole-original, conjecture, Lean or measurement credit
is claimed.
\end{document}
